% ============================================================
% CHAPTER 8: EXPERIMENTAL RESULTS AND EVALUATION
% ============================================================

\chapter{EXPERIMENTAL RESULTS AND EVALUATION}
\label{chap:results}

This chapter presents the quantitative evaluation of the QoS Scheduler system. The primary goal is to validate the effectiveness of the Token Bucket and WFQ algorithms in managing bandwidth and reducing latency under heavily congested network conditions without requiring root access.

\section{Experimental Setup}
\label{sec:experimental_setup}

The experiments were conducted using an Android device acting as a Wi-Fi Hotspot, connecting multiple client devices. Network traffic was generated and measured using standard benchmarking tools.

\begin{itemize}
    \item \textbf{Host Device:} Android 14 (Google Pixel 7), Snapdragon 8 Gen 2, 8GB RAM.
    \item \textbf{Client Devices:} Windows 11 PC (Traffic Generator) and an Android tablet.
    \item \textbf{Measurement Tools:} \texttt{iPerf3} for throughput injection, \texttt{ping} for latency/jitter analysis, and \texttt{Wireshark} for packet-level capture.
    \item \textbf{Network Baseline:} The upstream cellular connection provided a baseline capacity of approximately 20 Mbps downlink and 10 Mbps uplink.
\end{itemize}

\section{Throughput Enforcement}
\label{sec:results_throughput}

The most critical test measures the system's ability to enforce bandwidth caps on \texttt{LOW} priority traffic while maintaining throughput for \texttt{HIGH} priority streams.

In this scenario, a background client initiated a massive TCP file transfer (classified as \texttt{LOW} priority). Simultaneously, the primary client initiated a critical video stream (classified as \texttt{HIGH} priority).

\begin{table}[H]
\centering
\caption{Throughput distribution under extreme congestion (Target Cap: 5 Mbps for LOW)}
\label{tab:throughput_results}
\begin{tabular}{@{}lccc@{}}
\toprule
\textbf{Traffic Class} & \textbf{Baseline (No QoS)} & \textbf{With QoS Enabled} & \textbf{Improvement} \\
\midrule
HIGH Priority & 3.2 Mbps & 14.8 Mbps & \textbf{+362\%} \\
LOW Priority & 16.5 Mbps & 4.9 Mbps & (Throttled) \\
Total Network & 19.7 Mbps & 19.7 Mbps & 0\% \\
\bottomrule
\end{tabular}
\end{table}

As shown in Table~\ref{tab:throughput_results}, without QoS, the aggressive TCP file transfer consumed 83\% of the available bandwidth, starving the video stream. Upon enabling the QoS Scheduler, the \texttt{LOW} priority traffic was strictly throttled to 4.9 Mbps (successfully enforcing the 5 Mbps Token Bucket cap). Consequently, the \texttt{HIGH} priority traffic achieved a 3.2$\times$ throughput advantage, stabilizing at 14.8 Mbps.

\section{Latency and Jitter Reduction}
\label{sec:results_latency}

Bufferbloat is a pervasive issue in mobile hotspots, leading to severe latency spikes during heavy uploads or downloads. The system's WFQ implementation is designed to mitigate this by isolating traffic queues.

We measured ICMP Echo Request (ping) latency to a regional server (base latency $\sim$25ms) while the network was saturated with UDP traffic.

\begin{table}[H]
\centering
\caption{Latency and Jitter measurements during network saturation}
\label{tab:latency_results}
\begin{tabular}{@{}lccc@{}}
\toprule
\textbf{Metric} & \textbf{Baseline (No QoS)} & \textbf{With QoS Enabled} & \textbf{Reduction} \\
\midrule
Average Latency & 245.3 ms & 38.6 ms & \textbf{84.3\%} \\
Peak Latency (Spike) & 890.0 ms & 65.2 ms & \textbf{92.6\%} \\
Jitter (Variance) & 112.5 ms & 8.4 ms & \textbf{92.5\%} \\
Packet Loss & 4.2 \% & 0.1 \% & \textbf{97.6\%} \\
\bottomrule
\end{tabular}
\end{table}

The results in Table~\ref{tab:latency_results} are striking. Without QoS, the FIFO queue of the default Android hotspot quickly filled up, leading to severe bufferbloat (average latency of 245ms with spikes nearly reaching 1 second). By actively shaping the traffic and dropping packets that exceed the \texttt{LOW} priority threshold, the QoS Scheduler maintained an average latency of 38.6ms. Jitter was virtually eliminated (reduced from 112.5ms to 8.4ms), ensuring a stable connection for real-time applications like VoIP and competitive gaming.

\section{System Overhead and Profiling}
\label{sec:results_overhead}

Operating a VPN service and performing Deep Packet Inspection (DPI) in user-space introduces processing overhead. We profiled the Android device using Android Studio Profiler during a continuous 50 Mbps throughput test.

\begin{itemize}
    \item \textbf{CPU Utilization:} The \texttt{QosVpnService} process consumed an average of 12.4\% CPU. The majority of the load was concentrated in the Kotlin Coroutine IO Dispatcher managing the TCP/UDP relays.
    \item \textbf{Memory Footprint:} The system maintained a highly stable memory footprint of $\sim$68 MB. The aggressive reuse of \texttt{ByteBuffer} instances through a memory pool effectively prevented Garbage Collection (GC) pauses.
    \item \textbf{Battery Impact:} Over a 60-minute continuous test, the application accounted for only 4.5\% of total system battery drain, well within acceptable bounds for a background utility application.
\end{itemize}

\section{Statistical Significance}
\label{sec:results_statistics}

To ensure the reliability of the findings, the throughput and latency experiments were repeated 30 times ($N=30$). An independent samples t-test was conducted comparing the Baseline and QoS-Enabled groups. 

The reduction in latency was highly statistically significant ($M_{baseline} = 245.3$, $SD_{baseline} = 45.2$; $M_{qos} = 38.6$, $SD_{qos} = 5.1$; $t(58) = 24.8$, $p < .001$). The Cohen's $d$ effect size was calculated at 6.4, indicating a massive practical effect.

\section{Chapter Summary}
\label{sec:ch8_summary}

The empirical evaluation conclusively demonstrates that the QoS Scheduler successfully shapes traffic at the application level. By strictly enforcing Token Bucket caps on background traffic, the system successfully eliminates bufferbloat, reducing latency by over 84\% and stabilizing jitter. Furthermore, the user-space implementation proved highly efficient, imposing negligible overhead on the host device's CPU and battery.
